\subsection{KZG takeaways}

\begin{enumerate}
\item
  \emph{Elliptic curves} are very useful in cryptography. Roughly
  speaking, they are sets of points (usually in \(\mathbb{F}_{p}^{2}\))
  that satisfy some group law/``addition.'' The BN254 curve is a good
  ``typical curve'' to keep in mind.
\item
  The \emph{discrete logarithm} assumption is a common ``hard problem
  assumption'' used in cryptography with different groups. Specifically,
  since elliptic curves are groups, the discrete logarithm assumption
  over elliptic curves is very often used.
\item
  \emph{Commitment schemes} are ways for one party to commit values to
  another. Elliptic curves enable \emph{Pedersen commitments}, a very
  useful example of a commitment scheme.
\item
  Specifically, \emph{polynomial commitment schemes} are commitments of
  polynomials that are small and easy to ``open'' (evaluate at different
  points). KZG is one of the main polynomial commitment schemes being
  used in cryptography, such as in PLONK (coming up).
\end{enumerate}
